\chapter{API Reference}
\label{chap:api_reference}

This chapter provides a structured reference for the main classes and
functions exposed by \actsim. It is intended for use as a lookup
reference rather than a tutorial; readers who are new to \actsim should
start with the chapters on the individual classes.

% =====================================================================
\section{\texttt{DistributionFitter}}
\label{api:distribution_fitter}

\subsection*{Constructor}

\begin{lstlisting}[style=pythonstyle]
DistributionFitter(data, distributions=None, metrics=None)
\end{lstlisting}

\begin{longtable}{p{0.25\textwidth}p{0.7\textwidth}}
\toprule
\textbf{Argument} & \textbf{Description} \\
\midrule
\endhead
\texttt{data}          & Array-like of observed values to fit. \\
\texttt{distributions} & Optional list of distribution names. Defaults to
                         all supported distributions. \\
\texttt{metrics}       & Optional list of goodness-of-fit metrics.
                         Defaults to \texttt{['aic', 'bic']}. \\
\bottomrule
\end{longtable}

\subsection*{Methods}

\begin{longtable}{p{0.32\textwidth}p{0.62\textwidth}}
\toprule
\textbf{Method} & \textbf{Description} \\
\midrule
\endhead
\texttt{fit()} &
    Fits all candidate distributions to the data and computes
    goodness-of-fit metrics. \\
\texttt{select\_best\_fit()} &
    Internal method called by \texttt{fit()}; selects the best fit under
    each requested metric. \\
\texttt{get\_best\_fit(metric)} &
    Returns the best-fit dictionary for the given metric. \\
\texttt{select\_distribution(name)} &
    Selects a previously fitted distribution by name as the active
    distribution. Requires \texttt{fit()} to have been called. \\
\texttt{get\_selected\_dist()} &
    Returns the currently selected distribution object. \\
\texttt{get\_selected\_params()} &
    Returns the parameters of the currently selected distribution. \\
\texttt{predict(x)} &
    Evaluates the PDF (or PMF) of the selected distribution at \texttt{x}. \\
\texttt{sample(size=1)} &
    Draws random samples from the selected distribution. \\
\texttt{sample\_mixed(zero\_prop=0, one\_prop=0, size=1)} &
    Draws samples with a given proportion of zeros and ones inserted. \\
\texttt{calculate\_statistics()} &
    Returns mean, standard deviation, and key percentiles for both the
    data and the predicted density. \\
\texttt{plot\_predictions(\dots)} &
    Plots a histogram of the data with the PDFs of the requested
    distributions overlaid. Argument:
    \texttt{distribution\_names=None}. \\
\texttt{summary()} &
    Returns a DataFrame summarising the fitted distributions. \\
\texttt{truncate\_data(i)} &
    Removes observations equal to \texttt{i} from the data. \\
\bottomrule
\end{longtable}

% =====================================================================
\section{\texttt{StochasticSimulator}}
\label{api:stochastic_simulator}

\subsection*{Constructor}

\begin{lstlisting}[style=pythonstyle]
StochasticSimulator(
    freq_dist, freq_params,
    sev_dist, sev_params,
    num_sim=10000,
    keep_all=False,
    seed=1,
    correlation=None,
    copula_type=None,
    theta=0,
)
\end{lstlisting}

\begin{longtable}{p{0.27\textwidth}p{0.68\textwidth}}
\toprule
\textbf{Argument} & \textbf{Description} \\
\midrule
\endhead
\texttt{freq\_dist}    & Name of the frequency distribution. \\
\texttt{freq\_params}  & Tuple of frequency distribution parameters. \\
\texttt{sev\_dist}     & Name of the severity distribution. \\
\texttt{sev\_params}   & Tuple of severity distribution parameters. \\
\texttt{num\_sim}      & Number of simulation periods. \\
\texttt{keep\_all}     & If \texttt{True}, retain individual event records. \\
\texttt{seed}          & Random seed. \\
\texttt{correlation}   & Optional correlation between $N$ and $X$. \\
\texttt{copula\_type}  & Optional copula family. \\
\texttt{theta}         & Copula dependence parameter. \\
\bottomrule
\end{longtable}

\subsection*{Methods}

\begin{longtable}{p{0.36\textwidth}p{0.58\textwidth}}
\toprule
\textbf{Method} & \textbf{Description} \\
\midrule
\endhead
\texttt{gen\_agg\_simulations()} &
    Runs the aggregate loss simulation. \\
\texttt{gen\_copula()} &
    Generates correlated uniforms from the configured copula family. \\
\texttt{gen\_multivariate\_corr\_simulations(\dots)} &
    Generates correlated marginal simulations across multiple lines of
    business. Arguments: \texttt{corr\_matrix\_file},
    \texttt{dist\_list\_file}, \texttt{gen\_marginal=False}. \\
\texttt{calc\_agg\_percentile(pct=95)} &
    Returns the requested percentile of aggregate losses. \\
\texttt{plot\_distribution(bins=None, log\_option=False)} &
    Plots a histogram of simulated aggregate losses. \\
\texttt{plot\_correlated\_variables()} &
    Plots simulated frequency--severity correlation. \\
\texttt{analyze\_results(quantiles=[0.1, 0.25, 0.5, 0.75, 0.9, 0.95, 0.99], **kwargs)} &
    Returns VaR, TVaR, OEP, and AEP at the specified quantiles. \\
\texttt{apply\_deductible\_and\_limit(\dots)} &
    Applies per-occurrence and aggregate deductibles and limits.
    Arguments: \texttt{per\_occurrence\_ded},
    \texttt{per\_occurrence\_limit}, \texttt{agg\_ded},
    \texttt{agg\_limit}. \\
\bottomrule
\end{longtable}

\subsection*{Properties}

\begin{longtable}{p{0.3\textwidth}p{0.65\textwidth}}
\toprule
\textbf{Property} & \textbf{Description} \\
\midrule
\endhead
\texttt{results} &
    Pandas Series of aggregate losses by simulation period. \\
\texttt{all\_simulations} &
    Pandas DataFrame of individual event records (when
    \texttt{keep\_all=True}). \\
\bottomrule
\end{longtable}

% =====================================================================
\section{\texttt{ClaimSimulator}}
\label{api:claim_simulator}

\subsection*{Constructor}

\begin{lstlisting}[style=pythonstyle]
ClaimSimulator(
    policies_df,
    random_seed=42,
    correlation=None,
    copula_type=None,
    copula_param=0,
)
\end{lstlisting}

\begin{longtable}{p{0.27\textwidth}p{0.68\textwidth}}
\toprule
\textbf{Argument} & \textbf{Description} \\
\midrule
\endhead
\texttt{policies\_df}   & Policy table; must contain the columns
                         \texttt{policy\_id}, \texttt{freq\_dist},
                         \texttt{freq\_params}, \texttt{sev\_dist},
                         \texttt{sev\_params}, \texttt{start\_date},
                         and \texttt{end\_date}. The \texttt{freq\_params}
                         and \texttt{sev\_params} entries must be tuples. \\
\texttt{random\_seed}   & Random seed. \\
\texttt{correlation}    & Optional frequency--severity correlation. \\
\texttt{copula\_type}   & Optional copula family. \\
\texttt{copula\_param}  & Copula dependence parameter. \\
\bottomrule
\end{longtable}

\subsection*{Methods}

\begin{longtable}{p{0.36\textwidth}p{0.58\textwidth}}
\toprule
\textbf{Method} & \textbf{Description} \\
\midrule
\endhead
\texttt{group\_policies()} &
    Groups policies by frequency/severity specifications. \\
\texttt{simulate\_claims()} &
    Simulates claim counts and severities for each policy group. \\
\texttt{simulate\_dates\_nhpp(lambda0=10, alpha=0.5, phase=0, T=1)} &
    Assigns occurrence dates using a non-homogeneous Poisson process. \\
\texttt{simulate\_claim\_development(\dots)} &
    Develops claims using stochastic LDFs. Arguments:
    \texttt{base\_LDFs}, \texttt{volatility=0.1},
    \texttt{cumulative\_factor=1.0}. \\
\texttt{save\_claim\_development(filepath)} &
    Saves the development DataFrame to CSV. \\
\bottomrule
\end{longtable}

\subsection*{Attributes}

\begin{longtable}{p{0.3\textwidth}p{0.65\textwidth}}
\toprule
\textbf{Attribute} & \textbf{Description} \\
\midrule
\endhead
\texttt{policies}            & Processed policy DataFrame. \\
\texttt{claim\_data}         & Simulated claim records. \\
\texttt{claim\_development}  & Long-format claim development data. \\
\bottomrule
\end{longtable}

% =====================================================================
\section{\texttt{Config} and \texttt{load\_config}}
\label{api:config}

\subsection*{\texttt{load\_config}}

\begin{lstlisting}[style=pythonstyle]
load_config(file_path=None)
\end{lstlisting}

Convenience function that returns a \texttt{Config} object. If
\texttt{file\_path} is \texttt{None}, the default \texttt{config.yaml}
shipped with the package is used.

\subsection*{\texttt{Config}}

\begin{longtable}{p{0.3\textwidth}p{0.65\textwidth}}
\toprule
\textbf{Method} & \textbf{Description} \\
\midrule
\endhead
\texttt{has\_key(key)} &
    Returns \texttt{True} if the key exists in the configuration. \\
\texttt{update(updates)} &
    Updates configuration values from a dictionary. \\
\texttt{clear()} &
    Removes all configuration values. \\
\texttt{reload()} &
    Re-reads the configuration file from disk. \\
\texttt{save()} &
    Writes the current configuration to file. \\
\bottomrule
\end{longtable}

% =====================================================================
\section{Utility Functions}
\label{api:utilities}

\subsection*{\texttt{timing\_decorator}}

A decorator used internally to measure the runtime of long-running
methods. It prints the elapsed time when the decorated function returns.
This decorator is intended for development use rather than as part of
the public API.

% =====================================================================
\section{Supported Distribution Names}
\label{api:dist_names}

The following table lists the distribution names accepted by both
\texttt{DistributionFitter} and \texttt{StochasticSimulator}. The
underlying implementations are provided by the \texttt{actstats}
package (\texttt{actstats.actuarial}).

\begin{longtable}{p{0.3\textwidth}p{0.65\textwidth}}
\toprule
\textbf{Name} & \textbf{Description} \\
\midrule
\endhead
\texttt{normal}            & Normal (Gaussian) distribution. \\
\texttt{logistic}          & Logistic distribution. \\
\texttt{exponential}       & Exponential distribution. \\
\texttt{gamma}             & Gamma distribution. \\
\texttt{beta}              & Beta distribution. \\
\texttt{lognormal}         & Lognormal distribution. \\
\texttt{weibull}           & Weibull distribution. \\
\texttt{pareto}            & Pareto distribution. \\
\texttt{uniform}           & Uniform distribution. \\
\texttt{poisson}           & Poisson distribution. \\
\texttt{negative binomial} & Negative binomial distribution. \\
\bottomrule
\end{longtable}

Parameter conventions follow those of \texttt{actstats}; see the
\texttt{actstats} documentation for the exact parametrisation of each
distribution.
